\chapter{群、循环群与最初的对称操作}
\label{chap:group-basics}

群论研究的并不是某一种特定的对象，而是“可逆操作如何复合”的共同结构。对物理学而言，旋转、反射、晶格平移、粒子置换与量子态上的酉变换看起来十分不同，但只要它们的复合满足相同的代数规律，就可以在同一套语言中讨论。

把一个操作做完再做另一个操作，得到的仍应是允许的操作；什么也不做要算其中一个操作；每一步还应当能够撤销。封闭性、结合律、单位元和逆元并不是为了凑出一张定义表，而是把这几句实验语言压缩成没有歧义的公理。循环群和正三角形的对称操作会让每条公理立即工作。

本章采用乘法记号表示一般群运算。只有在加法群中才改用“$+$”、零元“$0$”和加法逆元“$-g$”。若 $A,B\subseteq G$，记
\[
AB\equiv\{ab\mid a\in A,\ b\in B\}.
\]
有限集合 $X$ 的元素个数记为 $\lvert X\rvert$。
\section{群、群元与循环结构}
\label{sec:group-axioms}


群的结构从一个二元运算开始。这个词第一次听起来很抽象，但它说的事情非常朴素：我们有一个集合，从中取两个对象，按照某条固定规则把它们合成第三个对象。例如两个整数可以相加，两个矩阵可以相乘，两个几何变换可以先后执行，两个置换也可以先做一个再做另一个。这些“合成规则”在数学上都属于二元运算。

这里最需要改变的一个直觉是：以后写 $ab$ 时，$a$ 与 $b$ 不一定是数，$ab$ 也不一定是普通乘法。它只表示“把 $a$ 与 $b$ 按当前规定的群运算复合”。如果群元是空间变换，我们约定 $ab$ 表示先做 $b$、再做 $a$；如果群元是矩阵，这正好与矩阵乘法作用在列向量上的顺序一致。因此，群论中的“乘法”更接近“操作复合”而不是初等算术中的乘法。

还有一个记号需要在这里读懂。写
\[
\mu:G\times G\longrightarrow G
\]
时，$G\times G$ 是所有有序对 $(a,b)$ 的集合；箭头表示 $\mu$ 是一个函数，它把每一对输入 $(a,b)$ 指派给唯一的输出。所谓“二元”，就是输入有两个位置；所谓“在 $G$ 上”，就是输入和输出都来自同一个集合。

\begin{definition}[二元运算]
设 $G$ 为非空集合。一个二元运算是映射
\[
\mu:G\times G\longrightarrow G,
\qquad
(a,b)\longmapsto \mu(a,b).
\]
若记 $\mu(a,b)=ab$，则映射值 $ab$ 必须仍属于 $G$；这就是运算的封闭性。
\label{def:binary-operation}
\end{definition}

\begin{definition}[群]
一个群是有序对 $(G,\cdot)$，其中 $G$ 为非空集合，$\cdot$ 为 $G$ 上的二元运算，并满足：
\begin{enumerate}[label=(\roman*),leftmargin=3em]
  \item \textbf{结合律：}对任意 $a,b,c\in G$，
  \[
  (ab)c=a(bc).
  \]
  \item \textbf{单位元存在：}存在 $e\in G$，使任意 $g\in G$ 满足
  \[
  eg=ge=g.
  \]
  \item \textbf{逆元存在：}对任意 $g\in G$，存在 $h\in G$，使
  \[
  gh=hg=e.
  \]
\end{enumerate}
若另外对任意 $a,b\in G$ 都有 $ab=ba$，则称 $G$ 为\textbf{交换群}或 \textbf{Abel 群}。
\label{def:group}
\end{definition}

封闭性已经包含在“二元运算 $G\times G\to G$”的定义中，因此不必把它与后面三条重复计数。单位元与逆元的唯一性也不应写入公理，因为它们可以由公理推出。

这三条公理各自解决一个不同问题。结合律保证连续做三个以上操作时不必记录括号，例如 $(ab)c=a(bc)$ 允许我们直接写成 $abc$；单位元 $e$ 描述“什么也不做”的操作；逆元 $g^{-1}$ 则描述“把 $g$ 的作用完全撤销”。群论之所以适合描述对称操作，正是因为物理上的可逆对称变换天然具有这三种性质。注意公理中\emph{没有}要求 $ab=ba$。绕不同轴的空间转动、一般矩阵以及置换通常都不交换，而非交换性正是后面许多结构出现的来源。

初学时可以把群公理理解为一个最小的“可逆操作系统”：任意两个允许的操作复合后仍然允许；复合顺序的括号无关；存在不改变任何东西的操作；每个操作都能被另一个操作撤销。只要这四层意思成立，我们就不再关心群元究竟是数字、矩阵、函数还是几何变换。

\begin{proposition}[单位元与逆元的唯一性]
群 $G$ 的单位元唯一；对每个 $g\in G$，其逆元唯一，记为 $g^{-1}$。
\label{prop:identity-inverse-unique}

\end{proposition}

\begin{proof}
若 $e$ 与 $e'$ 都是单位元，则
\[
e=e e'=e'.
\]
故单位元唯一。若 $h,k$ 都是 $g$ 的逆元，则
\[
h=h e=h(gk)=(hg)k=e k=k.
\]
故逆元唯一。
\end{proof}

\begin{proposition}[消去律与逆元公式]
对任意 $a,b,c\in G$，有
\[
ab=ac\Longrightarrow b=c,
\qquad
ba=ca\Longrightarrow b=c.
\]
并且
\[
(ab)^{-1}=b^{-1}a^{-1},
\qquad
(g^{-1})^{-1}=g.
\]
\label{prop:cancellation-inverse}

\end{proposition}

\begin{proof}
由 $ab=ac$ 左乘 $a^{-1}$，利用结合律得到
\[
a^{-1}(ab)=a^{-1}(ac)
\Longrightarrow
(a^{-1}a)b=(a^{-1}a)c
\Longrightarrow b=c.
\]
右消去律同理。另一方面，
\[
(ab)(b^{-1}a^{-1})
=a(bb^{-1})a^{-1}=e,
\]
且
\[
(b^{-1}a^{-1})(ab)
=b^{-1}(a^{-1}a)b=e,
\]
故由逆元唯一性得到 $(ab)^{-1}=b^{-1}a^{-1}$。
\end{proof}

\begin{definition}[群的阶]
若 $G$ 为有限集，则称 $\lvert G\rvert$ 为群 $G$ 的\textbf{阶}。若 $G$ 含无限多个元素，则称 $G$ 为无限群。
\label{def:group-order}
\end{definition}

\medskip
\noindent\emph{三个基本例子。}
整数集 $\ZZ$ 在加法下构成 Abel 群，单位元为 $0$，$n$ 的逆元为 $-n$。非零实数集 $\RR^{\times}=\RR\setminus\{0\}$ 在通常乘法下构成 Abel 群。$n\times n$ 可逆复矩阵在矩阵乘法下构成一般线性群 $\operatorname{GL}(n,\CC)$；当 $n\ge 2$ 时它一般不是 Abel 群。
\par\medskip

\medskip
\noindent\emph{置换群。}
设 $X=\{1,2,\dots,n\}$。所有双射 $\sigma:X\to X$ 在映射复合下构成群，记为对称群 $S_n=\operatorname{Sym}(X)$。其阶为
\[
\lvert S_n\rvert=n!.
\]
乘法约定为函数复合：$(\sigma\tau)(x)=\sigma(\tau(x))$，即右边的置换先作用。
\par\medskip

在继续引入新定义之前，可以先把群公理真正用于一次计算。考虑三个元素的集合
\[
C_3=\{e,a,a^2\},\qquad a^3=e.
\]
这里 $a$ 可以想成正三角形绕中心逆时针转 $120^\circ$。于是 $a^2$ 是转 $240^\circ$，再做一次 $a$ 就回到单位操作。它的乘法完全由指数相加决定：
\[
a\,a=a^2,\qquad a\,a^2=a^3=e,\qquad a^2a^2=a^4=a.
\]
若把结果排成乘法表，
\[
\begin{array}{c|ccc}
\cdot&e&a&a^2\\\hline
e&e&a&a^2\\
a&a&a^2&e\\
a^2&a^2&e&a
\end{array}
\]
每一行、每一列都恰好出现全部三个元素一次。这不是额外公理，而是消去律的结果：若同一行里 $ax=ay$，左消去就迫使 $x=y$。这个很小的例子把单位元、逆元、结合律和消去律都放到了同一张表里。

有了群公理，下一步自然是问：如果不断重复同一个操作 $g$，会产生哪些不同结果？例如把正三角形每次转 $120^\circ$，第三次就回到原位；而在整数加法群中不断加 $1$，永远不会回到 $0$。这两种行为分别对应“有限阶”和“无限阶”，也是循环子群概念的来源。

对 $g\in G$，定义
\[
g^0=e,
\qquad
g^n=\underbrace{gg\cdots g}_{n\text{ 个}},\quad n>0,
\qquad
g^{-n}=(g^{-1})^n.
\]
结合律保证了幂的括号位置无关，并有熟悉的指数律
\[
g^m g^n=g^{m+n},
\qquad
(g^m)^n=g^{mn},
\qquad m,n\in\ZZ.
\]
注意一般不能把 $(gh)^n$ 写成 $g^nh^n$；这需要 $gh=hg$。

\begin{definition}[生成子群与循环群]
由元素 $g\in G$ 生成的循环子群定义为
\[
\left\langle g\right\rangle=\{g^n\mid n\in\ZZ\}.
\]
若存在 $g\in G$ 使 $G=\left\langle g\right\rangle$，则称 $G$ 为循环群，$g$ 为 $G$ 的一个生成元。
\label{def:cyclic-group}
\end{definition}

\begin{proposition}[循环群必为 Abel 群]
若 $G=\left\langle g\right\rangle$，则 $G$ 为 Abel 群。

\end{proposition}

\begin{proof}
任取 $g^m,g^n\in G$，有
\[
g^m g^n=g^{m+n}=g^{n+m}=g^n g^m.
\]
\end{proof}

\begin{definition}[群元的阶]
对 $g\in G$，若存在最小正整数 $n$ 使 $g^n=e$，则称 $n$ 为 $g$ 的阶，记为
\[
\operatorname{ord}(g)=n.
\]
若不存在这样的正整数，则记 $\operatorname{ord}(g)=\infty$。
\label{def:element-order}
\end{definition}

\begin{proposition}[群元的阶与循环子群]
对任意 $g\in G$，
\[
\operatorname{ord}(g)=\lvert\left\langle g\right\rangle\rvert,
\]
其中等式允许两边同时为无穷。
\label{prop:element-order-cyclic}

\end{proposition}

\begin{proof}
若 $\operatorname{ord}(g)=n<\infty$，则 $e,g,\dots,g^{n-1}$ 两两不同；否则 $g^r=g^s$（$0\le r<s<n$）会给出 $g^{s-r}=e$，与 $n$ 的最小性矛盾。又因 $g^n=e$，任意 $g^k$ 都可按 $k=qn+r$ 化为 $g^r$，故
\[
\left\langle g\right\rangle=\{e,g,\dots,g^{n-1}\}.
\]
若 $\operatorname{ord}(g)=\infty$，则所有整数次幂两两不同。
\end{proof}

\begin{example}[$S_3$ 的基本结构]
采用循环记号，
\[
S_3=\{e,(12),(13),(23),(123),(132)\}.
\]
三个换位的阶均为 $2$，两个三循环的阶均为 $3$。正三角形的全部平面对称（或在三维中保持三角形集合不变的六个适当刚体操作）给出一个与 $S_3$ 同构的六阶群；物理教材中常记作 $D_3$。本讲义在抽象计算中优先使用 $S_3$，以避免不同文献对 $D_n$ 阶数约定的歧义。
\end{example}

$D_3$ 是阶最小的非 Abel 群，也是物理上最常见的点群之一（$C_{3v}$ 与它同构）。正因为它足够小，每一个后续概念——子群、共轭类、正规子群、商群、不可约表示——都可以在 $D_3$ 上完整算一遍；又因为它非交换，$D_3$ 已经能展示 Abel 群看不到的结构。下面先用几何语言把它的六个元素和乘法表固定下来，以后各节只需对这张表做不同的读法。

令 $r$ 为绕正三角形中心逆时针转 $120^\circ$，$s$ 为过顶点 $1$ 与对边中点所在直线的反射。则 $r$ 生成三阶循环子群 $\left\langle r\right\rangle=\{e,r,r^2\}$，$s$ 是二阶元。反射与旋转不交换：先转 $120^\circ$ 再反射，与先反射再转 $120^\circ$，把顶点送到不同位置。把这一非交换关系写成代数，就是
\begin{equation}
srs=r^{-1}=r^2,
\label{eq:d3-srs-relation}
\end{equation}
它也等价于 $sr=r^2s$ 与 $rs=sr^2$。于是全部六个元素可以统一写成
\[
D_3=\{e,\ r,\ r^2,\ s,\ sr,\ sr^2\},
\]
其中 $\{e,r,r^2\}$ 是纯转动，$\{s,sr,sr^2\}$ 是反射。\figref{fig:d3-symmetry} 把 $r$ 与 $s$ 的几何作用画在同一张图上：$r$ 是顶点 $1\to2\to3\to1$ 的循环置换，$s$ 固定顶点 $1$ 而交换 $2\leftrightarrow3$。

\begin{figure}[htbp]
\centering
\begin{tikzpicture}[line width=0.75pt,>=Stealth,every node/.style={font=\small}]
\coordinate (A1) at (90:1.6);
\coordinate (B1) at (210:1.6);
\coordinate (C1) at (330:1.6);
\draw (A1) -- (B1) -- (C1) -- cycle;
\fill (A1) circle (1.6pt) node[above=7pt] {$1$};
\fill (B1) circle (1.6pt) node[below left=13pt] {$2$};
\fill (C1) circle (1.6pt) node[below right=13pt] {$3$};
\draw[blue,mid arrow] (A1) to[bend left=28] (B1);
\draw[blue,mid arrow] (B1) to[bend left=28] (C1);
\draw[blue,mid arrow] (C1) to[bend left=28] (A1);
\node[blue] at (0,0.05) {$r$};
\node at (0,-2.3) {rotation $r$};
\begin{scope}[xshift=4.6cm]
\coordinate (A2) at (90:1.6);
\coordinate (B2) at (210:1.6);
\coordinate (C2) at (330:1.6);
\draw (A2) -- (B2) -- (C2) -- cycle;
\fill (A2) circle (1.6pt) node[above=7pt] {$1$};
\fill (B2) circle (1.6pt) node[below left=13pt] {$2$};
\fill (C2) circle (1.6pt) node[below right=13pt] {$3$};
\draw[dashed,red] (A2) -- (0,-1.9);
\draw[red,mid arrow] (B2) to[bend right=30] (C2);
\draw[red,mid arrow] (C2) to[bend right=30] (B2);
\node[red] at (0.34,0.20) {$s$};
\node at (0,-2.3) {reflection $s$};
\end{scope}
\end{tikzpicture}
\caption{正三角形的两个生成元：$r$ 为逆时针转 $120^\circ$，$s$ 为过顶点 $1$ 的反射。蓝箭头表示 $r$ 把每个顶点送到下一个，红色虚线为 $s$ 的反射轴。}
\label{fig:d3-symmetry}
\end{figure}

有了元素和关系 $r^3=s^2=e$、$sr=r^2s$，乘法表可以逐格算出而不依赖几何。例如 $r\cdot s=rs=sr^2$，$(sr)\cdot(sr)=s(rs)r=s(sr^2)r=r^3=e$，最后一式说明每个反射的阶都是 $2$。完整乘法表（行元素在左、列元素在右，乘积 $ab$ 表示先做 $b$ 再做 $a$）为
\begin{equation}
\begin{array}{c|cccccc}
D_3 & e & r & r^2 & s & sr & sr^2 \\ \hline
e     & e & r & r^2 & s & sr & sr^2 \\
r     & r & r^2 & e & sr^2 & s & sr \\
r^2   & r^2 & e & r & sr & sr^2 & s \\
s     & s & sr & sr^2 & e & r & r^2 \\
sr    & sr & sr^2 & s & r^2 & e & r \\
sr^2  & sr^2 & s & sr & r & r^2 & e
\end{array}
\label{eq:d3-multiplication-table}
\end{equation}
表中左上 $3\times3$ 块是 $\left\langle r\right\rangle\cong C_3$ 的乘法表，右下 $3\times3$ 块同样回到 $\left\langle r\right\rangle$，而两个非对角块把反射与转动互相转换——这正是非交换性的直接面貌。$D_3$ 可以由 $r,s$ 和三条关系完全生成：
\[
D_3=\bigl\langle r,\,s\,\bigm|\, r^3=s^2=e,\ srs=r^{-1}\bigr\rangle,
\]
这是群的\emph{生成元与关系}（presentation）写法。后面遇到 $D_n$、点群和空间群时，都优先用生成元与关系来描述，而不必列出全部元素。

\begin{derivation}[从生成元关系逐格算出乘法表]
乘法表中每一格都是行元素乘列元素，只需反复使用 $r^3=e$、$s^2=e$ 和关键关系 $sr=r^2s$（等价地 $rs=sr^2$）把乘积化回 $\{e,r,r^2,s,sr,sr^2\}$ 之一。

转动块。
$r\cdot r=r^2$，$r\cdot r^2=r^3=e$，$r^2\cdot r^2=r^4=r$，这些只用 $r^3=e$。

转动乘反射。
$r\cdot s=rs=sr^2$（直接用 $rs=sr^2$）。
$r\cdot sr=(rs)r=(sr^2)r=sr^3=s$。
$r\cdot sr^2=(rs)r^2=(sr^2)r^2=sr^4=sr$。
$r^2\cdot s=r(rs)=r(sr^2)=(rs)r^2=(sr^2)r^2=sr$。
$r^2\cdot sr=(r^2s)r=(sr)r=sr^2$。
$r^2\cdot sr^2=(r^2s)r^2=(sr)r^2=sr^3=s$。

反射乘转动。
$s\cdot r=sr$，$s\cdot r^2=sr^2$，按定义。
$sr\cdot r=sr^2$，$sr\cdot r^2=sr^3=s$。
$sr^2\cdot r=sr^3=s$，$sr^2\cdot r^2=sr^4=sr$。

反射乘反射。
$s\cdot s=s^2=e$。
$s\cdot sr=s^2r=r$，$s\cdot sr^2=s^2r^2=r^2$。
$sr\cdot s=s(rs)=s(sr^2)=s^2r^2=r^2$。
$sr\cdot sr=s(rs)r=s(sr^2)r=s^2r^3=e$，故 $(sr)^2=e$，即 $sr$ 是二阶元。
$sr\cdot sr^2=s(rs)r^2=s(sr^2)r^2=s^2r^4=r$。
$sr^2\cdot s=s(r^2s)=s(sr)=s^2r=r$（其中 $r^2s=sr$ 已在上面得到）。
$sr^2\cdot sr=s(r^2s)r=s(sr)r=s^2r^2=r^2$。
$sr^2\cdot sr^2=s(r^2s)r^2=s(sr)r^2=s^2r^3=e$，故 $(sr^2)^2=e$。

每个反射的平方都是 $e$，这与几何一致：反射做两次回到原位。非对角块中转动与反射的互相转换，则来自 $sr=r^2s$ 这一条非交换关系。
\end{derivation}

这个乘法表还可压缩为半直积 $D_3=C_3\rtimes C_2$。正规子群 $C_3=\langle r\rangle$ 与子群 $C_2=\langle s\rangle$ 满足 $C_3\cap C_2=\{e\}$ 且 $C_3C_2=D_3$，而 $srs^{-1}=r^{-1}$ 给出 $C_2$ 对 $C_3$ 的作用。于是每个元素唯一写成 $r^is^j$，乘法为
\begin{equation*}
(r^i s^j)(r^{i'} s^{j'})=r^{i+(-1)^j i'}s^{j+j'}
\end{equation*}
其中 $(-1)^j$ 体现了共轭作用 $s^jr^{i'}s^{-j}=r^{(-1)^ji'}$。这也是 $sr\ne rs$ 的代数根源。

在三个顶点上的作用给出同构 $D_3\cong S_3$：可取 $r\mapsto(1\,2\,3)$、$s\mapsto(2\,3)$。同一结构也由展示
\begin{equation*}
D_3=\langle r,s\mid r^3=e,\ s^2=e,\ srs^{-1}=r^{-1}\rangle .
\end{equation*}
概括。第三条关系正是半直积作用；它与前两条关系一起确定全部六个元素及其乘法。几何上，$r$ 的 $120^\circ$ 转动和 $s$ 的镜面反射给出 $D_3\hookrightarrow \mathrm O(2)$ 的忠实二维表示。

\figref{fig:d3-cayley-graph} 是 $D_3$ 的 Cayley 图：每个节点是一个群元，蓝色有向边表示右乘 $r$，红色无向边表示右乘 $s$（因 $s^2=e$ 故无方向）。Cayley 图把乘法表压缩成一张可读的图：从任一节点出发，沿蓝边走三步回到原处（$r^3=e$），沿红边走两步回到原处（$s^2=e$），而"蓝—红"与"红—蓝"到达的节点不同（$sr\neq rs$）。

\begin{figure}[htbp]
\centering
\begin{tikzpicture}[line width=0.75pt,>=Stealth,every node/.style={font=\small}]
% Inner triangle = <r> = {e,r,r^2}; outer triangle = {s,sr,sr^2}.
% s-edges are radial spokes, so no edge crossings.
\node[circle,draw,fill=white,inner sep=2pt,minimum size=11mm] (e)  at (90:1.7)  {$e$};
\node[circle,draw,fill=white,inner sep=2pt,minimum size=11mm] (r)  at (210:1.7) {$r$};
\node[circle,draw,fill=white,inner sep=2pt,minimum size=11mm] (r2) at (330:1.7) {$r^2$};
\node[circle,draw,fill=white,inner sep=2pt,minimum size=11mm] (s)   at (90:3.3)  {$s$};
\node[circle,draw,fill=white,inner sep=2pt,minimum size=11mm] (sr2) at (210:3.3) {$sr^2$};
\node[circle,draw,fill=white,inner sep=2pt,minimum size=11mm] (sr)  at (330:3.3) {$sr$};
% inner 3-cycle, right x r
\draw[lecturearrow,blue] (e) -- (r);
\draw[lecturearrow,blue] (r) -- (r2);
\draw[lecturearrow,blue] (r2) -- (e);
% outer 3-cycle, right x r
\draw[lecturearrow,blue] (s) -- (sr);
\draw[lecturearrow,blue] (sr) -- (sr2);
\draw[lecturearrow,blue] (sr2) -- (s);
% radial s-edges (undirected, dashed red)
\draw[red,line width=0.75pt,dashed] (e) -- (s);
\draw[red,line width=0.75pt,dashed] (r) -- (sr2);
\draw[red,line width=0.75pt,dashed] (r2) -- (sr);
\node[blue] at (2.15,0.55) {$r$};
\node[red] at (0.55,2.45) {$s$};
\end{tikzpicture}
\caption{$D_3$ 的 Cayley 图。蓝色有向边为右乘 $r$（转动），红色虚边为右乘 $s$（反射）。内三角为 $\left\langle r\right\rangle$，外三角为三个反射。}
\label{fig:d3-cayley-graph}
\end{figure}
